\chapter{Conclusion}
\label{ch:conclusion}

\section{Summary of Contributions}
\label{sec:ch8-summary}

This thesis has studied fixed-budget informed-corrector timing for masked
diffusion language models. The central question is when schedule-gain
optimisation is reducible to marginal signal ranking, and when it requires
interaction-aware or search-based scheduling. The thesis addresses this through
a model- and corrector-agnostic mathematical framework and an empirical case
study on ProSeCo-OWT.

The five formal contributions are as follows. First, the thesis formulates
fixed-budget corrector timing as a schedule-gain problem, separated from
remasking, token selection, predictor scheduling, and corrector-kernel design.
Second, Theorem~A establishes a marginal proxy regret bound: under approximate
additivity and proxy calibration, top-$B$ ranking of separable step-scores
achieves bounded regret against the exhaustive schedule oracle. Third,
Theorem~B and Theorem~B$'$ provide a pairwise-surrogate regret framework with
a finite-pool high-probability form that can be tested experimentally without
leakage. Fourth, Diagnostic Framework~C operationalises regime classification.
Fifth, the ProSeCo-OWT case study provides prior baseline evidence that tested
separable rankers do not recover the best-of-$N$ Monte Carlo oracle headroom,
while coordinate-descent greedy (CD-G) and beam-search greedy (BS-AG) search
procedures recover 74--84\% and 49--64\% of that headroom at $B\in\{2,3,4\}$.

\section{Limitations}
\label{sec:ch8-limitations}

Several important limitations qualify the conclusions. The empirical claims
rest on prior results pending Phase~0 reproducibility re-confirmation. Until
the pre-flight assertion suite (PF1--PF8) passes on the real ProSeCo-OWT
backend and the K=3 smoke test reproduces the qualitative pattern, these
statements are planning evidence rather than final thesis evidence.

The empirical scope is single-backbone. A bounded LLaDA-SFT probe did not
establish transferable MC-oracle headroom at the tested resolution. Whether
the regime characteristics observed on ProSeCo-OWT generalise to other masked
diffusion backends remains an open question.

The search procedures CD-G and BS-AG use true schedule-gain feedback and are
structural existence results. They demonstrate that better-than-uniform
schedules exist, but do not constitute a deployable scheduler. A deployable
Level-3 feature-conditioned scheduler requires a predictor trained on training
seeds and evaluated on held-out seeds; this has not yet been constructed.

All empirical conclusions are conditional on the chosen quality functional
$F$ (negative GPT-2 negative log-likelihood in the ProSeCo-OWT case study).

\section{Future Work}
\label{sec:ch8-future-work}

Immediate next steps are methodological. Phase~0 re-confirmation is the
blocking prerequisite: PF1--PF8 must pass with the real ProSeCo-OWT backend,
then the K=3 smoke test must reproduce the qualitative pattern, and only then
should K=30 replication and new experiments proceed. Phase~1 interaction
diagnostics should estimate pairwise interactions $\xi_{t,t'}$ at Levels~1
and~2 to classify ProSeCo-OWT under the Diagnostic Framework~C protocol.
If pairwise structure is positive and schedule-level validation of $Q(S)$
against $G(S)$ succeeds on a no-leakage candidate pool, Phase~2 pairwise
scheduling can proceed.

Longer-term directions include Level-3 feature-conditioned schedulers,
extension to additional backbones, and investigation of whether non-separable
quality-guided refinement (such as PRISM used as a full refinement backend
rather than a separable step-score) provides headroom beyond the tested
separable ranker class.
